\section{Chapter 5 Refresher: The Transformer --- Attention as the Core Computational Primitive --- Attention, Shapes, and Complexity}

\subsection*{Setting the Scene}
Attention replaces recurrence's sequential bottleneck with explicit pairwise comparisons implemented as batched matrix multiply---linear algebra becomes architecture instead of notation garnish.

\paragraph{Step A --- bilinear compatibility.}
\textbf{Problem.} Position \(t\) must aggregate information from other positions weighted by relevance.
\textbf{Insight.} Dot products score compatibility between transformed representations (queries vs keys).
\textbf{Lineage.} Bahdanau-style attention (2015); scaled dot-product formulation crystallised in Vaswani et al.\ (2017).
\textbf{Mental model.} Think ``matching embeddings,'' not mysticism---everything reduces to tensor contractions.

\paragraph{Step B --- softmax as differentiable selector.}
\textbf{Problem.} Hard argmax blocks gradients; we need smooth convex combinations.
\textbf{Insight.} Softmax (Chapter~1 tools) turns logits into mixing weights---Preface Step~10 intuition specialised row-wise.
\textbf{Lineage.} Statistical mechanics exponential tilts $\to$ neural attention gates.
\textbf{Mental model.} Temperature inherited from decoding chapters scales sharpness identically in spirit.

\paragraph{Step C --- masks encode causal laws.}
\textbf{Problem.} Autoregressive generation forbids peeking rightward.
\textbf{Insight.} Additive \(-\infty\) constraints zero softmax mass---enforce equality constraints algebraically before nonlinear squash.
\textbf{Lineage.} Decoder masking folklore codified in Transformer stacks.
\textbf{Mental model.} Same trick handles padding tokens---different semantics, identical mathematics.

\paragraph{Step D --- quadratic surface area.}
\textbf{Problem.} Pairwise attention allocates \(T\times T\) logits per layer.
\textbf{Insight.} Long contexts explode memory/compute---motivates FlashAttention-style engineering outside this refresher's scope but inside systems reasoning.
\textbf{Lineage.} Dao et al.\ IO-aware kernels (2022 era).
\textbf{Mental model.} Preface Step~6 variance arguments resurface when chunking randomises attention pathways.

If shapes align, attention is honest matrix cookbookery---when shapes drift, bugs hide silently.

\subsection*{Notation Introduced in This Chapter}
\begin{itemize}[leftmargin=1.5em]
  \item $T$: sequence length.
  \item $d_{\text{model}}$: model hidden width.
  \item $H$: number of attention heads.
  \item $d_k$: per-head key/query dimension.
  \item $Q,K,V$: query/key/value matrices.
  \item $M$: additive mask matrix applied before softmax.
\end{itemize}

\subsection*{What You Need To Recall Precisely}
\begin{itemize}[leftmargin=1.5em]
  \item \textbf{Scaled attention:}
  \[
    \mathrm{Attn}(Q,K,V)=\mathrm{softmax}\!\left(\frac{QK^\top+M}{\sqrt{d_k}}\right)V.
  \]
  Read row-wise: row $t$ of $QK^\top$ lists dot-product similarities between position-$t$ query vectors and every key vector; softmax turns those logits into nonnegative weights summing to $1$; multiplying those weights into $V$ mixes value vectors---attention outputs are convex combinations of values attended-to.

  \textbf{Why divide by $\sqrt{d_k}$?}
  If components of queries and keys are $\mathcal{O}(1)$ with roughly independent variation across $d_k$ dimensions, dot-product magnitudes grow \(\Theta(\sqrt{d_k})\) in expectation (length scales like Euclidean norm). Without scaling, softmax logits become sharply peaked as width grows, driving gradients toward sparse argmax-like behaviour and harming trainability. Rescaling keeps typical logits \(\mathcal{O}(1)\) across widths.

  \item \textbf{Mask semantics:} causal mask forbids looking right; padding mask blocks non-content locations.
  \item \textbf{Head decomposition:} independent projections create multiple interaction subspaces.
  \item \textbf{Residual and normalization:} residual streams pass representations forward as \(x\leftarrow x+\mathrm{Block}(x)\), letting blocks learn corrections around identity rather than full maps from scratch. Normalisation (LayerNorm/RMSNorm inside Transformer blocks) recentres/rescales activations per token or position so dot products do not drift orders of magnitude mid-training---pairs with residual paths to control gradient scale through depth.
  \item \textbf{Complexity source:} score matrix is $T\times T$, so interaction cost grows with sequence length.
\end{itemize}

\subsection*{How These Results Build on Each Other}
\begin{enumerate}[leftmargin=1.5em]
  \item Project input vectors into $Q,K,V$.
  \item Form pairwise compatibility scores via $QK^\top$.
  \item Apply structural constraints with mask $M$.
  \item Normalize row-wise with softmax.
  \item Aggregate values and return contextualized representations.
\end{enumerate}

\subsection*{Reminder Glossary (Term by Term)}
\begin{itemize}[leftmargin=1.5em]
  \item \textbf{Query:} representation of what current position is asking for.
  \item \textbf{Key:} representation of what each position offers.
  \item \textbf{Value:} content payload mixed by attention weights.
  \item \textbf{Mask:} hard structural restriction on allowable dependencies.
  \item \textbf{Head:} one parallel attention subspace.
  \item \textbf{KV cache:} stored keys/values for efficient autoregressive decoding.
\end{itemize}

\subsection*{Worked Mini Example (Shape Sanity)}
Let $T=4$, $d_{\text{model}}=8$, heads $H=2$, so $d_k=4$.
\[
Q,K,V \in \mathbb{R}^{4\times 4}\ \text{per head},\quad QK^\top \in \mathbb{R}^{4\times 4}.
\]
After softmax over rows, multiply by $V$ to return $\mathbb{R}^{4\times 4}$ per head, concatenate to $\mathbb{R}^{4\times 8}$.

\subsection*{Worked Example (Reasoning Chain with Causal Mask)}
For token position $t=3$ in a length-4 causal sequence, allowed keys are positions $\{1,2,3\}$ only.  
In row 3 of $QK^\top$, mask sets column 4 score to $-\infty$ before softmax.  
Result: attention weight at $(3,4)$ is exactly 0.  
Interpretation: causality is enforced at score level, not corrected later at output level.

\subsection*{Intuition Checkpoints}
\begin{itemize}[leftmargin=1.5em]
  \item ``Masking'' is a constraint on attention scores, not a post-hoc output edit.
  \item Quadratic attention cost refers to score interactions, not all block components equally.
  \item KV cache changes decoding-time behavior, not pretraining-time complexity class.
\end{itemize}

\subsection*{Further Refresh}
\begin{itemize}[leftmargin=1.5em]
  \item Shape and vector geometry reminders: see Chapter 4 refresher.
  \item Softmax stability reminders: see Chapter 1 refresher.
  \item \textbf{Vaswani et al. (2017)} --- architectural equations and design rationale.
  \item \textbf{Annotated Transformer} --- equation-by-equation implementation mapping.
\end{itemize}
